\section{DrainSinkhorn}
\label{sec:method}

\subsection{Problem and objective}

Consider $W$ independent entropic optimal-transport (EOT) problems executed by
one batched Sinkhorn backend. Standard fixed-width batching applies the same
per-problem Sinkhorn map in parallel but retains all $W$ lanes until the
slowest problem satisfies the configured stopping rule. When completion depths
differ, already-finished lanes therefore remain in later batched kernels. Our
objective is to remove this padding without changing any problem's EOT
objective, Sinkhorn update, stopping tolerance, two-marginal release rule, or
downstream consumer.

\subsection{Method overview}

DrainSinkhorn begins from the same fixed-width batch and repeatedly performs a
batched update, a preliminary one-marginal screen, the configured two-marginal
verification for screen-positive lanes, and physical compaction of every lane
that passes. For problem $t$, let $n_t$ be the batched-iteration number of the
first configured verification at which the problem passes. DrainSinkhorn then
writes its output and removes all of its per-problem state. Batched update $k$
therefore executes at survivor width
\[
A_k=\left|\{t:n_t\geq k\}\right|.
\]
The verification interval is the number of batched iterations between
consecutive verifications.
Logical masking records completion at width $W$; dynamic-compaction execution runs
  subsequent kernels at survivor width $A_k$. A fixed-width/dynamic-compaction comparison
measures the net effect of release and dynamic width reduction. LM/DC holds the
  release decisions fixed to isolate dynamic width reduction. The experiments
  report width and grouping interventions separately.

\subsection{Standard batched Sinkhorn execution}

Let $F_t$ denote one full Sinkhorn cycle for problem $t$, including that
problem's marginals, kernel operator, and numerical representation. Standard
batching adds a leading problem-batch dimension but no cross-problem reduction.
In ideal
arithmetic, its defining algebraic property is
\begin{equation}
\widetilde F\!\left(\operatorname{pack}(z_1,\ldots,z_W)\right)
=\operatorname{pack}\!\left(F_1(z_1),\ldots,F_W(z_W)\right).
\label{eq:pack-commutes}
\end{equation}
Thus the batched kernel preserves the EOT objective and per-instance Sinkhorn
map while changing their execution layout. After a verification, DrainSinkhorn
filters the problem-batch dimension and applies
Equation~\eqref{eq:pack-commutes} to the survivors.

For one problem, a Sinkhorn cycle performs
\begin{equation}
u^{k+1}=a\oslash(Kv^k),\qquad
v^{k+1}=b\oslash(K^\top u^{k+1}),
\label{eq:sinkhorn-cycle}
\end{equation}
where $\oslash$ denotes elementwise division. Immediately after the second
leg, exact arithmetic gives
\begin{equation}
v^{k+1}\odot(K^\top u^{k+1})=b.
\label{eq:updated-marginal}
\end{equation}
The column marginal is current by construction; the row marginal
$u^{k+1}\odot(Kv^{k+1})$ can remain outside tolerance. DrainSinkhorn evaluates
that row residual for all active problems in one batched operation and sends
only screen-positive problems to the two-marginal verifier.

\begin{figure*}[t]
\centering
\setlength{\fboxsep}{6pt}
\fbox{\parbox{.13\textwidth}{\centering
\textbf{Batched update}\\[-1pt]
$u\leftarrow a/(Kv)$\\
$v\leftarrow b/(K^\top u)$}}
\hspace{2pt}$\rightarrow$\hspace{2pt}
\fbox{\parbox{.13\textwidth}{\centering
\textbf{One-marginal screen}\\[-1pt]
batch row residuals\\
updated column identity}}
\hspace{2pt}$\rightarrow$\hspace{2pt}
\fbox{\parbox{.13\textwidth}{\centering
\textbf{verification}\\[-1pt]
current row and column\\
stopping tolerance}}
\hspace{2pt}$\rightarrow$\hspace{2pt}
\fbox{\parbox{.13\textwidth}{\centering
\textbf{Release and compact}\\[-1pt]
emit passing problems\\
compact all OT state}}
\hspace{2pt}$\rightarrow$\hspace{2pt}
\fbox{\parbox{.11\textwidth}{\centering
\textbf{Narrower next kernel}\\[-1pt]
$W_k\to W_{k+1}$}}
\caption{DrainSinkhorn's execution loop. The preliminary screen using the current
marginals avoids many full plan evaluations. Problems that pass the
two-marginal verifier leave every
per-problem tensor; the next batched update processes only the survivors.}
\label{fig:active-loop}
\end{figure*}

\subsection{Screen, verify, and compact}

Let $\widehat r^{\mathrm{row}}_{t,k}$ be the batched screen value. A problem is
screen-positive when
$\widehat r^{\mathrm{row}}_{t,k}\leq\tau_{\mathrm{screen}}$; otherwise it
continues without materializing the more expensive two-marginal verification. For a
screen-positive problem, the implementation recomputes both marginals and applies
the configured stopping rule
\begin{equation}
r_{t,k}=\max\!\left\{
\|\pi_{t,k}\ones-a_t\|_1,
\|\pi_{t,k}^\top\ones-b_t\|_1
\right\}\leq\tau_{\mathrm{stop}}.
\label{eq:release-contract}
\end{equation}
A rejected problem remains active for a later verification. A false-positive
screen invokes the two-marginal verifier unnecessarily, whereas a false-negative
screen delays removal. We measure both effects by replaying the screen outcomes
and evaluating every replayed candidate with the verifier in
\cref{sec:experiments}.

For each passing problem, DrainSinkhorn writes the output and compacts every
per-problem object: source and target descriptors, marginals, dual potentials,
centered shifts, log weights, and scratch buffers. The next batched update sees
only surviving problems. The logical-mask
control uses the same update, preliminary screen, two-marginal verifier,
tolerance, and observed
completion times. It freezes a passing problem's potentials but retains that
problem in the batch tensor, so subsequent kernels still execute at the
original width. Logical masking versus dynamic compaction therefore isolates
the work removed by reducing the batch width.

Equation~\eqref{eq:pack-commutes} and
Equation~\eqref{eq:updated-marginal} define the ideal-arithmetic transformation.
Fixed-work parity, achieved residuals, near-threshold stress tests, and
application endpoints characterize its implementation behavior.

\paragraph{Stopping-rule-preserving release.}
A problem is removed only after passing the same two-marginal verifier used by
the selected backend. The preliminary screen using the current marginals decides
only whether to invoke that verification; it does not authorize removal. DrainSinkhorn
therefore preserves the backend's configured release authority while changing
how unfinished problems are
scheduled and stored. The timed packed Triton path applies this rule in the precision
supported by that backend. Independent FP64 replays serve as evaluation references
for finite-precision decision behavior.

\subsection{Execution modes and fallback}

The implementation exposes three batch execution modes. Fixed-width batching
updates every problem until a common final boundary. Logical masking records
per-problem completion and
freezes completed states without narrowing the kernel. Dynamic compaction
removes passing problems. Width determines whether a batched kernel is
efficient, and the numbers of iterations needed by the problems in the batched
window determine how much width dynamic compaction can remove. Formal
campaigns fix execution mode, width, verification interval, and
fallback before timing. Formal campaigns use fail-closed convergence handling:
iteration-limit and non-finite-state events raise an error. The packed Triton
implementation optionally hands a dynamic-compaction tail to the matched
single-problem solver when survivor width falls below a preconfigured
\texttt{min\_packed\_width}. The formal MetroPT-3 and Packer19 configurations
set this threshold to one.

When a complete problem fits on one accelerator, different windows are sent to
independent complete-device workers, giving each device an independent update
path and avoiding per-update collectives. Row sharding remains a capacity path
for a single oversized problem; the main scale-out path assigns complete
problems to independent devices.
